% Source: https://artofproblemsolving.com/community/c3h1250108_hm_geometry
% Stable identity: aops-topic-1250108-q01
Place $A=(0,0)$, $O=(2,0)$, $P=(2,1)$, $S=(0,1)$.

\begin{asy}
usepackage("amsmath");
size(4.5cm);
defaultpen(fontsize(10pt));
pair A=(0,0),O=(2,0),P=(2,1),S=(0,1),E=(2-sqrt(3),1);
draw(A--O--P--S--cycle,linewidth(1)); draw(A--E,blue); draw(E--O,red);
label("$A$",A,SW);label("$O$",O,SE);label("$P$",P,NE);label("$S$",S,NW);label("$E$",E,N);
label("$2$",(1,0),S); label("$1$",(2,0.5),E); draw(arc(E,0.25,180,degrees(A-E)));
label("\mbox{}",(min(currentpicture,user=true)+max(currentpicture,user=true))/2,black); // keeps the text after the figure black when LaTeX typesets the labels inline
\end{asy}

In right triangle $ASE$ (right angle at $S$), $\tan\angle AES=\dfrac{AS}{SE}$, so $SE=\dfrac{1}{\tan75^\circ}=\tan15^\circ=2-\sqrt3$. Thus $E=\left(2-\sqrt3,1\right)$ and
\[EO=\sqrt{\left(\sqrt3\right)^2+1^2}=\boxed{2}.\]
